\section[Cap\'{i}tulo 2: Aprendizaje Pr\'{a}ctico de Ciberseguridad]{\centering Cap\'{i}tulo 2: Aprendizaje Pr\'{a}ctico de Ciberseguridad}

El aprendizaje pr\'{a}ctico es necesario porque la ciberseguridad no se limita a saber conceptos, sino a usarlos en momentos donde la informaci\'{o}n puede ser incompleta o ambigua. Un estudiante puede conocer qu\'{e} es el \textit{phishing}, pero aun as\'{i} fallar si no sabe revisar dominios, remitentes o se\~{n}ales de urgencia falsa (Prince, 2004).

Este cap\'{i}tulo explica por qu\'{e} el proyecto se apoya en actividades interactivas. La propuesta no busca reemplazar la clase te\'{o}rica, sino complementarla con ejercicios donde el estudiante tome decisiones, observe consecuencias y refuerce su criterio de respuesta (Kolb, 1984).

\subsection{Limitaciones del enfoque te\'{o}rico}

La ense\~{n}anza expositiva permite introducir conceptos, pero no siempre alcanza para desarrollar habilidades de an\'{a}lisis y respuesta. En seguridad inform\'{a}tica, muchas tareas requieren pr\'{a}ctica: revisar un correo, interpretar una alerta, decidir si se escala un caso o contener un equipo comprometido (Prince, 2004).

En grupos numerosos, el docente puede tener dificultades para acompa\~{n}ar a cada estudiante durante ejercicios pr\'{a}cticos. Esto reduce oportunidades de retroalimentaci\'{o}n y puede dejar una brecha entre lo que el estudiante sabe explicar y lo que realmente puede resolver frente a un caso (Biggs \& Tang, 2011).

\subsection{Aprendizaje activo}

El aprendizaje activo plantea que el estudiante participa mediante an\'{a}lisis, resoluci\'{o}n de problemas y toma de decisiones. Este enfoque encaja con ciberseguridad porque obliga a usar evidencias, evaluar alternativas y justificar acciones frente a una situaci\'{o}n concreta (Prince, 2004).

\begin{center}
\begin{tabular}{|p{0.30\textwidth}|p{0.55\textwidth}|}
\hline
\textbf{Estrategia} & \textbf{Aplicaci\'{o}n en el videojuego} \\
\hline
An\'{a}lisis de casos & Correos, llamadas y alertas con informaci\'{o}n incompleta o sospechosa (Kolb, 1984). \\
\hline
Toma de decisiones & Acciones como aprobar, escalar, reportar o contener un incidente (Prince, 2004). \\
\hline
Aprendizaje por error & Consecuencias mostradas en reportes de turno y registros del sistema (Hattie \& Timperley, 2007). \\
\hline
Progresi\'{o}n gradual & Casos simples al inicio y escenarios combinados en niveles posteriores (Biggs \& Tang, 2011). \\
\hline
\end{tabular}
\end{center}

Desde esta perspectiva, el videojuego funciona como un espacio de pr\'{a}ctica controlada. El estudiante puede ensayar respuestas y equivocarse sin generar da\~{n}o real, lo cual resulta \'{u}til para temas donde el error en un entorno laboral puede tener consecuencias serias (Kolb, 1984).

\subsection{Retroalimentaci\'{o}n en el aprendizaje}

La retroalimentaci\'{o}n ayuda a reducir la distancia entre el desempe\~{n}o actual y el desempe\~{n}o esperado. Para que sea \'{u}til, debe entregar informaci\'{o}n sobre la acci\'{o}n realizada y orientar una mejora, no solo indicar si algo estuvo bien o mal (Hattie \& Timperley, 2007).

En el proyecto, la retroalimentaci\'{o}n se plantea mediante consecuencias dentro del escenario. Por ejemplo, si el jugador aprueba un correo sospechoso, el reporte posterior puede indicar que un equipo inici\'{o} actividad an\'{o}mala. As\'{i} se mantiene el tono de simulaci\'{o}n y, al mismo tiempo, el estudiante entiende el efecto de su decisi\'{o}n (Hattie \& Timperley, 2007).

\subsection{S\'{i}ntesis del cap\'{i}tulo}

El aprendizaje pr\'{a}ctico justifica el uso de escenarios interactivos en el proyecto. Para formar habilidades de ciberseguridad, el estudiante necesita observar, decidir y revisar consecuencias. El videojuego se plantea como una forma de apoyar ese proceso sin abandonar el contenido acad\'{e}mico.
